\chapter{Dyskusja i~wnioski}
\label{ch:dyskusja-wnioski}

\section{Odpowiedzi na pytania badawcze}
\label{sec:odpowiedz-pytania}

Pierwsze pytanie dotyczyło tego, jak zaprojektować pipeline wyliczający współczynniki biegu z~pojedynczego nagrania jedną kamerą i~jaką osiąga dokładność. Praca pokazała, że trójetapowy pipeline --- estymacja pozy przez MediaPipe, klasyfikacja faz siecią BiLSTM i~obliczenia geometryczne --- pozwala wyliczyć wszystkie dwanaście współczynników z~jednego nagrania (rozdział~\ref{ch:materialy-metody}). Dokładność oceniłem względem zakresów z~literatury, a~dla jednego nagrania także względem niezależnego pomiaru z~zegarka i opaski sportowej (rozdział~\ref{ch:wrazliwosc}). Kadencja zgodziła się z~pomiarem referencyjnym w~granicach 1\%, natomiast czas kontaktu z~podłożem był systematycznie zaniżony o~około 14\%. 

Drugie pytanie dotyczyło wyboru klasyfikatora faz i~jego wpływu na precyzję współczynników czasowych. Model sekwencyjny, uwzględniający kontekst sąsiednich klatek, przewyższył klasyfikator analizujący każdą klatkę niezależnie (rozdział~\ref{ch:klasyfikator}). Inżynieria cech poprawiła wyniki lasu losowego, ale większy zysk przyniosło dopiero wprowadzenie kontekstu czasowego w~sieci BiLSTM. Co istotne, dokładność klasyfikacji przekłada się nierównomiernie na precyzję poszczególnych metryk. Współczynniki oparte na zliczaniu zdarzeń, takie jak kadencja, okazały się odporne, podczas gdy współczynniki oparte na pomiarze czasu trwania faz, takie jak czas kontaktu z~podłożem, są wrażliwe na błędy granic faz. Wniosek ten wykazałem najpierw rachunkowo, a~następnie potwierdziłem niezależnym pomiarem referencyjnym.

Trzecie pytanie dotyczyło wrażliwości współczynników na warunki nagrania i~możliwości automatycznego wykrywania przypadków niewiarygodnych. Wykazałem, że współczynniki przestrzenne --- zwłaszcza wzorzec lądowania stopy --- są wrażliwe na perspektywę kamery, a~współczynniki czasowe na częstotliwość klatek (rozdział~\ref{ch:wrazliwosc}). Najbardziej zawodne przypadki, w~których wartość wykracza poza zakres fizjologicznie możliwy, udało się wykrywać automatycznie, bez pomiaru referencyjnego --- mechanizm ten zadziałał poprawnie nawet na nagraniu w~wysokiej rozdzielczości i~orientacji poziomej, gdzie problem perspektywy nie był oczywisty.

\section{Wkład pracy}
\label{sec:wklad-pracy}

Wkład pracy odpowiada lukom w~literaturze zidentyfikowanym w~rozdziale~\ref{ch:sota}. Po pierwsze, praca dotyczy biegu sportowego, a~nie chodu klinicznego, który dominuje w~dotychczasowych badaniach --- bieg ma odmienne fazy i~inne wartości referencyjne. Po drugie, zamiast proponować pojedynczy model, systematycznie porównałem cztery konfiguracje klasyfikatora na świadomie małym, dobrze opisanym zbiorze. Po trzecie, zbadałem wpływ korekty proporcji obrazu jako kroku preprocessingu dla nagrań o~nietypowej orientacji, czego wcześniejsze prace nie raportowały. Po czwarte, uporządkowałem wartości referencyjne pochodzące z literatury i zaimplementowałem je w postaci działającego silnika reguł rekomendacji. Każda reguła zawiera przypisane źródło oraz poziom ważności komunikatu.
 Po piąte, zaproponowałem i~wdrożyłem mechanizm automatycznego wykrywania niskiej wiarygodności pomiaru bez dostępu do danych referencyjnych.

Dodatkowym, niezaplanowanym pierwotnie wkładem jest walidacja względem zegarka sportowego. Jest to jedyny w~pracy punkt odniesienia pochodzący z~zewnętrznego urządzenia, a~nie z~zakresów literaturowych, i~to on dostarczył ilościowego dowodu na różną odporność metryk opartych na zliczaniu i~na pomiarze czasu trwania.

\section{Ograniczenia}
\label{sec:ograniczenia-pracy}

\begin{itemize}
    \item \textbf{Niewielki zbiór danych.} Piętnaście nagrań pochodzących od~kilkunastu biegaczy to bardzo niewielki zbiór danych. Ogranicza to możliwość uogólniania wyników na całą populację biegaczy.
    \item \textbf{Brak pełnego pomiaru referencyjnego 3D.} Nie dysponowałem systemem motion capture, więc nie mogę podać dokładności w~jednostkach fizycznych dla większości nagrań. 
\item \textbf{Systematyczne zaniżenie czasu kontaktu.} Porównanie z pomiarem referencyjnym pokazało, że pipeline zaniża czas kontaktu z podłożem o około 14\%. Prawdopodobną przyczyną jest sposób wyznaczania granic między fazą kontaktu i lotu oraz różnica między metodami pomiaru: Garmin korzysta z czujników ruchu, natomiast pipeline wyznacza fazy na podstawie obrazu.
\\
    \item \textbf{Ograniczenia obrazu dwuwymiarowego.} Współczynniki kątowe wyliczane są w~płaszczyźnie obrazu, więc ujęcie inne niż ściśle z boku zniekształca ich wartości.
    \item \textbf{Ogólny charakter wartości referencyjnych.} Progi reguł pochodzą z~literatury dla zdrowych biegaczy i~mogą nie odpowiadać indywidualnej sytuacji konkretnej osoby.
\item \textbf{Drobne różnice między uruchomieniami MediaPipe.} Przy ponownym przetwarzaniu tego samego nagrania MediaPipe może zwrócić minimalnie inne pozycje keypointów. Zwykle są to bardzo małe różnice, ale mogą wpływać na wynik, jeśli dana wartość znajduje się blisko progu reguły diagnostycznej.

\end{itemize}

\section{Warunki wiarygodnego działania systemu}
\label{sec:implikacje}

Z~analizy wynikają jasne wskazania, kiedy systemowi można ufać. Przy standardowym ujęciu z~boku (orientacja pozioma, kamera prostopadła do toru biegu) i~częstotliwości co najmniej 15~klatek na sekundę wiarygodne są zarówno współczynniki czasowe, jak i~wskaźniki symetrii. Przy ujęciu pionowym lub ukośnym wzorzec lądowania stopy staje się niewiarygodny, o~czym system ostrzega automatycznie. Przy niskiej częstotliwości klatek metryki czasowe tracą precyzję. System odpowiada na inną potrzebę niż laboratoryjne systemy referencyjne: udostępnia obiektywną analizę techniki biegu każdemu, kto dysponuje zwykłą kamerą --- bez kosztownego sprzętu, laboratorium i~obecności specjalisty. Jego realną siłą jest połączenie tej dostępności z~automatycznym sygnalizowaniem, którym wartościom można zaufać, a~które wymagają ostrożności --- cecha, której tradycyjne narzędzia pomiarowe nie oferują.

\section{Kierunki dalszych badań}
\label{sec:praca-przyszla}

Pracę można rozwijać w~kilku kierunkach.

\begin{itemize}
    \item \textbf{Większy i~bardziej zróżnicowany zbiór danych.} Rozmiar zbioru był głównym ograniczeniem pracy. Zebranie nagrań od~znacznie większej liczby biegaczy --- o~różnej technice, prędkości i~budowie ciała --- pozwoliłoby rzetelniej ocenić generalizację, ograniczyć przeuczenie klasyfikatora i~przejść na podział danych na poziomie biegaczy.
    \\
    \item \textbf{Dokładniejsze źródło keypointów.} MediaPipe wprowadza szum, który ogranicza precyzję współczynników, zwłaszcza kątowych. Warto zbadać dokładniejsze metody estymacji pozy, a~na etapie badawczym pozyskać keypointy z~systemu motion capture jako materiał o~wyższej jakości --- posłużyłyby one zarówno do zmierzenia, ile precyzji tracimy przez szum estymacji, jak i~jako lepszy punkt odniesienia przy dalszym rozwoju modelu.
    \item \textbf{Rekalibracja granic faz.} Wykryte w~porównaniu z~pomiarem Garmin systematyczne zaniżenie czasu kontaktu o~około 14\% wskazuje konkretny kierunek poprawy --- dostrojenie sposobu wyznaczania granic między fazą kontaktu a~fazą lotu, tak aby lepiej odpowiadały rzeczywistym momentom kontaktu i~oderwania stopy.
    \item \textbf{Przeliczanie wyników na jednostki fizyczne.} Wprowadzenie kalibracji skali obrazu (np.~na podstawie znanego wzrostu biegacza) pozwoliłoby raportować współczynniki takie jak oscylacja pionowa czy długość kroku w~centymetrach i~metrach zamiast w~wartościach znormalizowanych.
\end{itemize}

Pipeline można też uzupełnić o~automatyczne ostrzeganie użytkownika o~zbyt niskiej częstotliwości klatek lub niewłaściwej orientacji nagrania już na etapie wczytywania wideo --- co wynika bezpośrednio z~przeprowadzonej analizy wrażliwości.

\section{Wnioski końcowe}
\label{sec:wnioski-koncowe}

Praca wykazała, że z~pojedynczego nagrania biegu zwykłą kamerą można, z~użyciem uczenia maszynowego do klasyfikacji faz, wyliczyć komplet współczynników biomechanicznych i~wygenerować na ich podstawie rekomendacje oparte na literaturze. Kluczowym wynikiem jest rozróżnienie współczynników odpornych od~wrażliwych na warunki nagrania oraz wykazanie, że te najbardziej zawodne można rozpoznać automatycznie. System nie dorównuje dokładnością laboratoryjnym systemom referencyjnym, ale oferuje to, czego one nie dają --- dostępność i~uczciwą, automatyczną sygnalizację granic własnej wiarygodności.
